%%%%%%%%%%  scan idx 230  |  folio 217  %%%%%%%%%%
\[
  C(\overline U,C^\infty(V,F^\bullet))
  \longrightarrow C(\overline U,C_2^r(V,F^\bullet))
\]
est un
quasi-isomorphisme. Quitte à remplacer $S$ par $\overline U$ (et $E^\bullet$ par le complexe induit au-dessus de $\overline U$), on est ramené à montrer que
\[
  C(S,C^\infty(V,F^\bullet))
  \longrightarrow C(S,C_2^r(V,F^\bullet))
\]
est un quasi-isomorphisme, $S$ étant supposé \emph{compact}. D'après (3.2.1.1) on a
\[
  C(S,C^\infty(V,F^\bullet))=
  \varprojlim_r C(S,C_2^r(V,F^\bullet))
\]
Pour chaque $i$, $C(S,C_2^r(V,F^i))$ est un espace de Banach (pour la topologie de la convergence uniforme), et les morphismes de transition
\[
  C(S,C_2^p(V,F^i))\longrightarrow C(S,C_2^q(V,F^i))
\]
sont, en vertu de (3.0.4) et (3.1.2), les applications linéaires continues d'image dense. D'autre part, appliquant aux quasi-isomorphismes (a') $\underline F_p^\bullet\longrightarrow\underline F_q^\bullet$ le foncteur $\Gamma'_S$, qui est exact puisque $\cO_S$ est mou, on voit que les morphismes de transition
\[
  C(S,C_2^p(V,F^\bullet))\longrightarrow C(S,C_2^q(V,F^\bullet))
\]
sont des quasi-isomorphismes. L'assertion (a) résulte alors de (EGA $0_{\mathrm{III}}$ 13.2.3, variante «~topologique~»). Ceci achève la démonstration de (3.2.1).

\textbf{3.2.2.} Si $S$ est un espace topologique, nous noterons $\Parf(S)$ la catégorie triangulée des complexes parfaits de $(\cO_S\otimes\mathbb C)$-Modules, et $K^\bullet(S)$ le «~groupe de Grothendieck~» correspondant\footnote{voir (I 6.3). Pour une revue générale des foncteurs $K$, voir Exp.~IV.}.

Soient $S$, $X$, $E^\bullet$ comme en (3.2.1)~; on appelle \emph{indice analytique} de $E^\bullet$ l'élément
\[
  \operatorname{ind.an}(E^\bullet)=C\ell(\mathbb{R}p_*(E^\bullet))\in K^\bullet(S).
\]
